\begin{ex}
Hình nào dưới đây không có cạnh thẳng và đỉnh nhọn? (1: Hình tròn; 2: Hình vuông; 3: Hình tam giác; 4: Hình chữ nhật).
\shortans{1}
\loigiai{
Hình tròn là hình khép kín bởi đường cong trơn đều, không có cạnh thẳng và không có đỉnh nhọn (ghi số 1).
\textbf{Đáp số:} $1$.
}
\end{ex}

\begin{ex}
Tổng số cạnh của một hình tam giác và một hình vuông là bao nhiêu?
\shortans{7}
\loigiai{
Hình tam giác có 3 cạnh, hình vuông có 4 cạnh.\nTổng số cạnh là: $3 + 4 = 7$ (cạnh).
\textbf{Đáp số:} $7$.
}
\end{ex}

\begin{ex}
Tổng số góc của hai hình tam giác là bao nhiêu?
\shortans{6}
\loigiai{
Mỗi hình tam giác có 3 góc.\nHai hình tam giác có: $3 + 3 = 6$ (góc).
\textbf{Đáp số:} $6$.
}
\end{ex}

\begin{ex}
Hiệu giữa số góc của một hình chữ nhật và một hình tam giác là bao nhiêu?
\shortans{1}
\loigiai{
Hình chữ nhật có 4 góc, hình tam giác có 3 góc.\nHiệu số góc là: $4 - 3 = 1$ (góc).
\textbf{Đáp số:} $1$.
}
\end{ex}

\begin{ex}
Một hình tam giác lớn được chia thành 3 hình tam giác nhỏ bằng các đường kẻ từ một đỉnh xuống đáy. Hỏi có tất cả bao nhiêu hình tam giác?
\shortans{6}
\loigiai{
Số hình tam giác đơn là 3 hình.\nSố hình tam giác ghép từ 2 hình đơn là 2 hình.\nSố hình tam giác ghép từ cả 3 hình đơn là 1 hình.\nTổng cộng có: $3 + 2 + 1 = 6$ (hình tam giác).
\textbf{Đáp số:} $6$.
}
\end{ex}

\begin{ex}
Một bức tranh gồm 4 hình chữ nhật và 5 hình tam giác. Hỏi trong tranh có bao nhiêu hình phẳng?
\shortans{9}
\loigiai{
Tổng số hình phẳng có trong bức tranh là: $4 + 5 = 9$ (hình).
\textbf{Đáp số:} $9$.
}
\end{ex}

\begin{ex}
Có 10 que tính, bạn Nam xếp được một hình chữ nhật có chiều dài 3 que tính, chiều rộng 2 que tính. Hỏi Nam còn thừa bao nhiêu que tính?
\shortans{0}
\loigiai{
Số que tính dùng để xếp hình chữ nhật là: $3 + 3 + 2 + 2 = 10$ (que tính).\nSố que tính còn thừa là: $10 - 10 = 0$ (que tính).
\textbf{Đáp số:} $0$.
}
\end{ex}

\begin{ex}
Bạn Lan có 6 hình tròn màu đỏ và 4 hình tròn màu vàng. Lan cất đi 3 hình tròn màu đỏ. Hỏi Lan còn lại bao nhiêu hình tròn cả hai màu?
\shortans{7}
\loigiai{
Số hình tròn đỏ Lan còn lại là: $6 - 3 = 3$ (hình tròn).\nTổng số hình tròn Lan còn lại là: $3 + 4 = 7$ (hình tròn).
\textbf{Đáp số:} $7$.
}
\end{ex}

\begin{ex}
Cần ít nhất bao nhiêu que tính để xếp được 2 hình vuông riêng biệt?
\shortans{8}
\loigiai{
Mỗi hình vuông cần 4 que tính. Hai hình vuông riêng biệt cần: $4 + 4 = 8$ (que tính).
\textbf{Đáp số:} $8$.
}
\end{ex}

\begin{ex}
Cần ít nhất bao nhiêu que tính để xếp được 3 hình tam giác riêng biệt?
\shortans{9}
\loigiai{
Mỗi hình tam giác cần 3 que tính. Ba hình tam giác riêng biệt cần: $3 + 3 + 3 = 9$ (que tính).
\textbf{Đáp số:} $9$.
}
\end{ex}

\begin{ex}
Dùng một nhát kéo thẳng cắt ngang qua một góc của hình tam giác. Hỏi mảnh giấy lớn còn lại có bao nhiêu cạnh?
\shortans{4}
\loigiai{
Cắt cụt 1 góc của hình tam giác tạo thêm 1 cạnh mới tại vết cắt. Số cạnh của hình lớn còn lại là: $3 - 1 + 2 = 4$ cạnh (trở thành hình tứ giác).
\textbf{Đáp số:} $4$.
}
\end{ex}

\begin{ex}
Bạn Hùng có 8 hình dán ngôi sao, bạn Dũng có 5 hình dán ngôi sao. Hỏi bạn Hùng có nhiều hơn bạn Dũng bao nhiêu hình dán ngôi sao?
\shortans{3}
\loigiai{
Số hình dán bạn Hùng nhiều hơn bạn Dũng là: $8 - 5 = 3$ (hình dán).
\textbf{Đáp số:} $3$.
}
\end{ex}

\begin{ex}
Để ghép thành một hình vuông lớn gồm 3 hàng và 3 cột, em cần bao nhiêu hình vuông nhỏ?
\shortans{9}
\loigiai{
Số hình vuông nhỏ cần dùng là: $3 \times 3 = 9$ (hình vuông nhỏ).
\textbf{Đáp số:} $9$.
}
\end{ex}

\begin{ex}
Một hình chữ nhật dài 6 cm, rộng 4 cm. Hiệu giữa chiều dài và chiều rộng của hình chữ nhật đó là bao nhiêu xăng-ti-mét?
\shortans{2}
\loigiai{
Hiệu giữa chiều dài và chiều rộng là: $6 - 4 = 2$ (cm).
\textbf{Đáp số:} $2$.
}
\end{ex}

\begin{ex}
Em có 7 que tính, em xếp được 1 hình vuông và còn thừa mấy que tính?
\shortans{3}
\loigiai{
Xếp 1 hình vuông dùng hết 4 que tính.\nSố que tính còn thừa là: $7 - 4 = 3$ (que tính).
\textbf{Đáp số:} $3$.
}
\end{ex}

\begin{ex}
Tổng số cạnh của 2 hình vuông là bao nhiêu cạnh?
\shortans{8}
\loigiai{
Mỗi hình vuông có 4 cạnh. Hai hình vuông có: $4 + 4 = 8$ (cạnh).
\textbf{Đáp số:} $8$.
}
\end{ex}

\begin{ex}
Tổng số cạnh của 3 hình tam giác là bao nhiêu cạnh?
\shortans{9}
\loigiai{
Mỗi hình tam giác có 3 cạnh. Ba hình tam giác có: $3 + 3 + 3 = 9$ (cạnh).
\textbf{Đáp số:} $9$.
}
\end{ex}

\begin{ex}
Một hình vuông có 4 góc. Ghép 2 hình vuông cạnh nhau tạo thành hình chữ nhật. Hỏi hình chữ nhật mới có bao nhiêu góc vuông?
\shortans{4}
\loigiai{
Bất kỳ hình chữ nhật nào cũng có đúng $4$ góc vuông.
\textbf{Đáp số:} $4$.
}
\end{ex}

\begin{ex}
Có 10 que tính, bé xếp được 2 hình tam giác có chung một cạnh (hết 5 que). Hỏi bé còn lại bao nhiêu que tính?
\shortans{5}
\loigiai{
Số que tính bé còn lại là: $10 - 5 = 5$ (que tính).
\textbf{Đáp số:} $5$.
}
\end{ex}

\begin{ex}
Trong các hình: hình tròn, hình tam giác, hình vuông, hình chữ nhật; có bao nhiêu hình có 4 cạnh?
\shortans{2}
\loigiai{
Các hình có 4 cạnh là hình vuông và hình chữ nhật. Vậy có $2$ hình.
\textbf{Đáp số:} $2$.
}
\end{ex}

\begin{ex}
Hình tam giác có ít hơn hình vuông bao nhiêu cạnh?
\shortans{1}
\loigiai{
Hình tam giác có 3 cạnh, hình vuông có 4 cạnh. Số cạnh ít hơn là: $4 - 3 = 1$ (cạnh).
\textbf{Đáp số:} $1$.
}
\end{ex}

\begin{ex}
Một bánh xe đạp quay tròn tạo thành vệt hình gì? (Hình tròn: 1, Hình vuông: 2, Hình tam giác: 3).
\shortans{1}
\loigiai{
Bánh xe đạp quay tròn tạo thành hình tròn (ghi số 1).
\textbf{Đáp số:} $1$.
}
\end{ex}

\begin{ex}
Một chiếc xúc xắc đồ chơi có các mặt là hình gì? (Hình tròn: 1, Hình tam giác: 2, Hình vuông: 3, Hình chữ nhật: 4).
\shortans{3}
\loigiai{
Khối xúc xắc có các mặt đều là hình vuông (ghi số 3).
\textbf{Đáp số:} $3$.
}
\end{ex}

\begin{ex}
Ghép 2 hình tam giác vuông bằng nhau có thể tạo thành hình nào? (Hình tam giác lớn hoặc hình vuông/chữ nhật: ghi số 1 nếu tạo được hình vuông, ghi 0 nếu không).
\shortans{1}
\loigiai{
Hai hình tam giác vuông cân bằng nhau ghép lại có thể tạo thành một hình vuông (ghi số 1).
\textbf{Đáp số:} $1$.
}
\end{ex}

\begin{ex}
Một lá cờ Tổ quốc Việt Nam có dạng hình chữ nhật và có 1 ngôi sao vàng. Hỏi lá cờ có bao nhiêu góc vuông?
\shortans{4}
\loigiai{
Lá cờ hình chữ nhật có đúng $4$ góc vuông.
\textbf{Đáp số:} $4$.
}
\end{ex}
